%% =============================================================================
%% PAPER 1 - Appendix
%% =============================================================================

\subsection{BTP Italia: Indexation Coefficient and Deflation Floors}
\label{app:btp_floor}

The BTP Italia indexes both coupons and principal to the Italian
consumer price index for blue- and white-collar worker households,
excluding tobacco (FOI ex tabacchi) published by Istat. At each coupon date $t$, the
indexation coefficient $\mathrm{CI}_{t}$ is the ratio of the relevant
FOI reference index to its value at the previous coupon date
$t-6\text{m}$, so that $\mathrm{CI}_{t}$ measures the cumulative
inflation over the six-month coupon period. The semiannual coupon and
the capital revaluation payment are computed as
\begin{equation}
C_{t} \;=\; \tfrac{1}{2}\, c\, N \,\widetilde{\mathrm{CI}}_{t},
\qquad
R_{t} \;=\; N\,\bigl(\widetilde{\mathrm{CI}}_{t}-1\bigr),
\label{eq:btp_italia_payments}
\end{equation}
where $c$ is the fixed real annual coupon rate set at issuance, $N$ is
the subscribed nominal capital, and $\widetilde{\mathrm{CI}}_{t}$ is
the \emph{modified} indexation coefficient that incorporates a
semiannual deflation floor:
\begin{equation}
\widetilde{\mathrm{CI}}_{t} \;=\; \max\!\bigl(\mathrm{CI}_{t},\, 1\bigr).
\label{eq:btp_italia_floor}
\end{equation}
If the FOI ex tabacchi falls between consecutive coupon dates, the
indexation coefficient is reset to one for that semester, the coupon
reverts to the fixed real rate $c/2 \cdot N$, and no capital
revaluation is paid; in the following semester, the base index is reset
to the higher of the previous semester's level and the highest reading
observed in earlier semesters. The principal is redeemed at par at
maturity ($N$), regardless of the cumulative inflation path. Two
features of this design are relevant for the arbitrage. First, the
semiannual reset of $\widetilde{\mathrm{CI}}_{t}$ to unity at every
coupon date means that the notional exposed to default risk is always
the original face value $N$, eliminating the path-dependent contamination
that affects conventional inflation-linked bonds such as the U.S. TIPS
or the euro-area BTP\euro i. Second, the deflation floor on the coupon
guarantees a strictly positive nominal payment in every period, which
is absent from instruments with floor only at maturity.


%% ---- A.2 Trade Construction Details ----
\subsection{Trade Construction Details}
\label{app:trade_details}

\subsubsection{Entry/Exit Rules and Thresholds}
\label{app:entry_exit}
Table~\ref{tab:entry_exit_rules} summarises entry and exit thresholds
for each strategy. All strategies require a minimum of six months to
maturity at entry and maintain at least one open trade after the first
signal. The CDS--Bond universe is restricted to single-name CDS whose
reference entity belongs to the current on-the-run iTraxx Main series;
names that exit the index are no longer traded. All CDS are centrally
cleared. The iTraxx skew strategy trades only on-the-run series across
Main, SnrFin, SubFin, and Xover.

\input{\tablepath/entry_exit_thresholds_article}

\subsubsection{Trade-Level Statistics}
\label{app:trade_stats}

Table~\ref{tab:trade_stats} reports trade-level statistics for each
strategy, including the number of trades, average duration, average
profit per trade, and the distribution of exit reasons.

\input{\tablepath/trade_statistics_article}

\subsubsection{Transaction Costs}
\label{app:fees}
Transaction costs are computed at the individual trade level and applied
as entry and exit fees proportional to DV01 (bonds) or duration (CDS).
Fee levels are regime-dependent: at each trade date, the prevailing level
of the iTraxx Main 5Y index determines which of three tiers applies
(Table~\ref{tab:fee_schedule}). Thresholds are set at 60 and 100~bps,
reflecting the empirical widening of bid-offer spreads in stressed markets.
For trades held to maturity, the exit fee is waived.

\input{\tablepath/fee_schedule_article}

%% ---- A.3 Signal-Weighted vs Equal-Weighted Comparison ----
\subsection{Signal-Weighted vs Equal-Weighted Comparison}
\label{app:sw_ew}

Table~\ref{tab:sw_ew_basic} compares the signal-weighted (SW,
\citealp{rebonato2021convexity}) and equal-weighted (EW,
\citealp{duarte2007risk}) indices. SW is adopted as the baseline in the
main text; EW results confirm robustness to the weighting scheme. The
equal-weighted counterpart of the manipulation-proof performance
measure (Table~\ref{tab:mppm}) is reported in
Table~\ref{tab:mppm_ew}; the volatility-managed results for both
weighting schemes are in Table~\ref{tab:moreira_muir}.

\input{\tablepath/SW_EW_basic_metrics_3strategies_article}

\input{\tablepath/MPPM_analysis_monthly_ew_article}

\input{\tablepath/moreira_muir_performance_monthly_2panel_article}


%% ---- Benchmark Regressions with US Factors (Section 4) ----
\subsection{Benchmark Regressions with US Factors}
\label{app:benchmark_us}
For completeness, the three benchmark regressions are re-estimated using the
original US factor proxies; the euro-area results in the main text
(Tables~\ref{tab:duarte}, \ref{tab:fung_hsieh}, and~\ref{tab:active_fi}) remain
the primary specification.

\input{\tablepath/Duarte_US_article_monthly}
\input{\tablepath/FungHsieh_US_article_monthly}
\input{\tablepath/ActiveFI_US_article_monthly}

%% ---- A.7 Subperiod and Rolling Window Analysis ----
\subsection{Subperiod and Rolling Window Analysis}
\label{app:subperiod}
 
This appendix reports sub-period splits, stress-regime decompositions,
and out-of-sample robustness for the benchmark factor models of
Section~\ref{sec:benchmarks}. Each table contains Panel~A (full sample,
first half, second half) and Panel~B (alpha by LOW/MEDIUM/HIGH stress
regime on the iTraxx Main 5Y level, thresholds at 60 and 100 bps).
 
 
%% ---- A.7.1 Duarte et al. (2007) ----
\subsubsection{Duarte et al.\ (2007)}
\label{app:subperiod_duarte}
Table~\ref{tab:subperiod_regime_duarte} reports sub-period splits and stress
regime alpha for the \citet{duarte2007risk} specification.
\input{\tablepath/subperiod_regime_duarte_monthly}
 

%% ---- A.7.2 Brooks et al. (2020) / Active FI ----
\subsubsection{Brooks et al.\ (2020) --- Active Fixed Income}
\label{app:subperiod_activefi}
Table~\ref{tab:subperiod_regime_activefi} reports the same analysis under the
\citet{brooks2020active} specification.
\input{\tablepath/subperiod_regime_activefi_monthly}
 
 
 
%% ---- A.7.3 Fung & Hsieh (2004) ----
\subsubsection{Fung \& Hsieh (2004)}
\label{app:subperiod_fh}
Table~\ref{tab:subperiod_regime_fh} reports results for the
\citet{fung2004hedge} specification.

 
\input{\tablepath/subperiod_regime_fh_monthly}
 

%% ---- A.7.4 Out-of-Sample Alpha: Rolling-Window Robustness ----
\subsubsection{Out-of-Sample Alpha: Rolling-Window Robustness}
\label{app:oos_rolling}
Table~\ref{tab:oos_benchmark_rolling} re-estimates the out-of-sample
alpha of Panel~C of Table~\ref{tab:alpha_synthesis} with the loadings
fitted on a rolling 60-month window, which relaxes the constant-loadings
assumption of the expanding scheme. The layout mirrors Panel~C, so the
two schemes can be compared cell by cell; the expanding estimates are
not reproduced here.

\input{\tablepath/oos_alpha_appendix_monthly}

%% ---- A.9 VIF Diagnostics ----
\subsection{VIF Diagnostics}
\label{app:vif}
 
Table~\ref{tab:vif_benchmarks} reports Variance
Inflation Factors for the EUR factor specifications used in the benchmark
regressions of Section~\ref{sec:benchmarks}. VIF is computed on the factor matrix for each
strategy's sample period; because the factors are identical across strategies,
differences in VIF reflect only the different estimation windows.
 
Two groups of regressors are highly collinear. In the Duarte
specification, the maturity-sorted bond portfolios and the bond index
returns carry elevated VIF (R5~$\approx$~52--55, R10~$\approx$~24--27,
R2 and RI~$\approx$~14--16), the expected consequence of the common
duration exposure of government and investment-grade bond returns. In
the Brooks et al.\ specification, the two global bond benchmarks are
near-duplicates (Global\_Term and Global\_Aggregate,
VIF~$\approx$~47--66). Collinearity among controls inflates the
standard errors of those loadings but leaves inference on the
intercept unaffected; all Fung--Hsieh factors have VIF below~4.
 
\input{\tablepath/VIF_combined_article_monthly}

%% ---- A.4 Candidate Factor Universe ----
\subsection{Candidate Factor List with Sources}
\label{app:factor_list}
Table~\ref{tab:factor_list} reports the full set of candidate risk
factors used in the principal-component and best-subset analyses of
Section~\ref{sec:factor_selection}, together with
their data sources, transformations, and economic category.
\input{\tablepath/factor_list}


%% ---- A.5 PCA Robustness ----
\subsection{PCA Robustness}
\label{app:pca_robustness}

\subsubsection{Robustness across the Number of Components}
\label{app:pca_wk}

Table~\ref{tab:pca_robustness} reports alpha and $\bar{R}^2$ across the
number of retained components $K$. Results are robust to the choice of $K$.

\input{\tablepath/PCA_robustness_alpha_R2_article}

\subsubsection{Subperiod and Regime Analysis}
\label{app:pca_subperiod}

Half-sample splits and regime-dependent alpha confirm that PCA-based
alpha is stable across sub-periods and increases during stress.

\input{\tablepath/PCA_subperiod_results_article}

\subsubsection{Conditional Alpha: Christopherson--Ferson--Glassman (PCA)}
\label{app:pca_conditional}

Table~\ref{tab:pca_cond_fs_contemporaneous_article} reports the
Christopherson--Ferson--Glassman conditional alpha using PCA factors:
because the model uses the full continuous variation in the stress
proxy rather than a discrete split, it establishes that the regime
effect of Table~\ref{tab:alpha_regime_pca_aen} does not hinge on the
choice of cutoffs; the cutoffs themselves are the 60 and 100
basis-point thresholds that govern the transaction-cost tiers
(Appendix~\ref{app:fees}).

\input{\tablepath/PCA_Cond_Alpha_FS_article_contemporaneous}




%% ---- A.6 Factor-Selection Robustness ----
\subsection{Factor-Selection Robustness}
\label{app:aen_robustness}

\subsubsection{Sensitivity to the Subset Cardinality}
\label{app:aen_sensitivity}

Table~\ref{tab:ml_ksens} re-estimates the best-subset selection as the
cardinality $k$ varies around the baseline. The selected sets are broadly
nested and the in-sample fit plateaus, so the conclusions do not hinge on
the choice of $k$.

\input{\tablepath/k_sensitivity}

\subsubsection{Factor Selection Across Methods}
\label{app:aen_selection_matrix}

Table~\ref{tab:aen_selection_names} reports the factor set selected for each
strategy by the best-subset primary selector and by two independent
cross-checks, the adaptive elastic net and post-double-selection. The
overlap across selectors is partial, as expected when highly correlated
candidates substitute for one another, but the economically central
exposures recur in every method, and the double-selection control set is a
strict subset of exactly these consensus exposures---the fixed-income
defensive factor for BTP~Italia, emerging-market debt and the industrial
bond return for the CDS--Bond basis---so the identified risks are not an
artifact of the cardinality constraint or of any single penalization
scheme.

\input{\tablepath/AEN_Selection_Names_article}


\subsubsection{Conditional Alpha: Detailed Results (Best-Subset)}
\label{app:aen_conditional}

Table~\ref{tab:conditional_alpha_fs} reports the full
Christopherson--Ferson--Glassman conditional alpha model with
best-subset-selected factors; as in the PCA case, the continuous conditioning
variable makes the regime effect independent of the choice of discrete
cutoffs.

\input{\tablepath/Cond_Alpha_FS_Thesis}


\subsubsection{Rolling Alpha}
\label{app:aen_rolling}

Figure~\ref{fig:app_aen_rolling_alpha}
plots rolling 36-month alpha with regime shading for the best-subset-selected factors.
Alpha spikes during stress episodes.

\begin{figure}[H]
\centering
\includegraphics[width=\textwidth]{\figpath/aen/aen_rolling_alpha.pdf}
\caption[Rolling Alpha with Regime Shading (Best-Subset)]{Rolling Alpha with Regime Shading (Best-Subset Factors).
Grey bands: 95\% confidence intervals.
Red shading: iTraxx Main 5Y $> 100$ bps. Rolling window: 36 months.}
\label{fig:app_aen_rolling_alpha}
\end{figure}



%% ---- A.8 Cross-Strategy Interdependencies: Additional Tests ----
\subsection{Cross-Strategy Interdependencies: Additional Tests}
\label{app:rq3_additional}
 
This appendix collects robustness checks and supplementary analyses for
Section~\ref{sec:interdependencies}. All tests use the same mispricing
magnitudes $m_{i,t}$ and regime definitions as the main text.
 
 
%% ---- A.8.1 DCC-GARCH ----
\subsubsection{DCC-GARCH}
\label{app:dcc}
 
We estimate bivariate DCC-GARCH(1,1) models on each pair
of $\Delta m$ series. Figure~\ref{fig:app_dcc} plots the time-varying
conditional correlations. Spikes in conditional correlation coincide with
stress episodes (COVID-19, ECB hiking cycle), corroborating the
regime-dependent results in Section~\ref{sec:interdependencies}.
 
\begin{figure}[H]
\centering
\includegraphics[width=\textwidth]{../results/rq3_duffie/figures/F1_dcc_garch.pdf}
\caption[Dynamic Conditional Correlations (DCC-GARCH)]{Dynamic Conditional Correlations (DCC-GARCH).
  Bivariate DCC(1,1) estimates for each pair of $\Delta m$.
  Shaded areas: HIGH stress regime (iTraxx Main $> 100$ bps).}
\label{fig:app_dcc}
\end{figure}
 
 

 
%% ---- A.8.3 Granger Causality: Lag Robustness ----
\subsubsection{Granger Causality: Lag Robustness}
\label{app:granger_multilag}

Table~\ref{tab:rq3_granger_multilag} reports Granger causality tests
at lag lengths $p = 1, 2, 3$. The main result --- iTraxx Skew
Granger-causes both the CDS--Bond basis and BTP~Italia --- is robust
across all specifications.

\input{\tablepath/RQ3_Granger_Multilag_article}

 